\documentclass{article}
\usepackage{amsmath,amssymb}
\begin{document}
\title{Nonvanishing of Dirichlet \(L\)-functions to the Right of One Half}
\author{}
\date{}
\maketitle

\section*{Main line}
The target is \(L(\rho,\chi)\ne0\) whenever \(\Re\rho>1/2\), except when \(\chi\) is principal and \(\rho=1\). Write \(R_M(y)\) for the one-sided prime error attached to a prime \(d\) and modulus \(M\).

Fix \(q,\chi,\rho\), and put
\[
 M=18480q,\qquad a=9240q-1,\qquad
 \varepsilon=\frac{\min(\Re\rho,1)-1/2}{51},\qquad
 \beta=\frac12+8\varepsilon .
\]
Then \(0<\varepsilon<1/100\) and \(1/2<\beta<\Re\rho\). Every prime \(d\equiv a\pmod M\) has the required arithmetic properties and gives parameters for Gaps 1 and 2. Those gaps yield, eventually,
\(R_M(x)\le x^{\beta-1}\). The allowed one-sided criterion makes the entire \(L\)-function of the character induced to modulus \(M\) nonzero at \(\rho\). Changing modulus and removing its finite Euler factors gives the target.

\section*{Gaps to prove independently}
\noindent\textbf{Gap 1.} Let \(M\) be a positive multiple of \(18480\), and let \(d\) be prime with \(M\mid2(d+1)\) but \(M\nmid d+1\). For any admissible geometric datum with scales \(x,x_1,b\), number of lifts \(m\), order \(A\), number of places \(f\), \(H>0\) thick places, and \(T>0\) top places, set
\[
 \mathcal B=fA(\log b)^3,\qquad
 \mathrm{rem}=AH\left(R_M(x)-R_M(x_1)+\frac TH R_1(x_1)\right).
\]
Prove \(m\,\mathrm{rem}\le(\log b)^2\mathcal B\). This is the covolume step: construct sufficiently many short lifts, bound their index, and compare the upper covolume bound with the product-formula lower bound. The common main height cancels, leaving the displayed inequality.

\medskip
\noindent\textbf{Gap 2.} For every \(M,d\) as in Gap 1 and every \(0<\varepsilon<1/100\), prove that every sufficiently large \(x\) has an admissible datum as above satisfying
\[
 m\ge\sqrt{x},\qquad \log b=x^\varepsilon,\qquad
 \frac fH\le\varphi(M)\log b,\qquad
 R_M(x_1)-\frac THR_1(x_1)\le x^{-10^6}.
\]
Establish the last bound by finding \(\log x_1\) in
\([10^6\log x,(10^6\log x)^2]\) where the indicated combination is small; the prime-error transforms supply the needed polar window argument. This gap enters after Gap 1: division by \(mAH\) gives
\[
 R_M(x)\le
 \frac{f(\log b)^5}{mH}+x^{-10^6}
 \le x^{8\varepsilon-1/2}
\]
for sufficiently large \(x\).

\section*{Allowed facts}
The one-sided prime-error criterion, finite Euler-factor behavior under change of modulus, and the elementary arithmetic of \(a\bmod M\) stated in the question.

\end{document}
